\section{A scalar limit from unbounded modulus}
\label{sec:scalar}

Assume that the macroscopic moduli $M(n)$ are unbounded. Crossing
duality first gives transverse weights of energy $O(M(n)^{-1})$ along
a subsequence. We turn those weights into continuous functions whose
level continua meet every path straddling the corresponding value.
An annular estimate then yields a nonconstant limit $u$ and a finite
measure controlling its squared oscillation. The approximating functions
retain vanishing local oscillation energy after each fixed group
transformation; later comparisons will use this property as well as the
limit itself. Constants are independent of the approximation index;
thresholds may depend on the fixed compact subset of the domain under
consideration.

\subsection{Record scales and crossing duality}

If $\limsup M(n)^{1/n}>1$, choose
$1<\lambda<\limsup M(n)^{1/n}$; otherwise put $\lambda=1$.
In both cases $M(n)/\lambda^n$ is unbounded. Retain increasing record
indices $N$ whose record values tend to infinity. Thus $M(N)\to\infty$ and
\begin{equation}
 M(N-l)\leq\lambda^{-l}M(N)\qquad(0\leq l\leq N).
 \label{eq:records}
\end{equation}
This inequality survives every further subsequence.

\begin{lemma}[Localization to a rectangle]
\label{lem:rectangle-localization}
There are coordinate rectangles
$\cQ=[0,1]^2\subset\Qhat=[0,1]\times[-1,2]\subset Z$ and a
subsequence of the record indices such that the left-to-right crossing
modulus of $\cQ$ is at least $cM(N)$, with $c>0$ independent of $N$.
\end{lemma}
\begin{proof}
Cover $Z$ by finitely many coordinate boxes $[-1,1]^2$, each lying in
a chart containing $[-8,8]^2$ of diameter less than $d_0$. A macroscopic
path meeting an inner box must leave its larger box. Follow it to its
first meeting with $\partial[-2,2]^2$. If it exits through $x=2$,
its segment after the last preceding meeting with $x=1$ crosses
$[1,2]\times[-2,2]$ from left to right. The other three exit sides
give analogous strips. Each strip extends by its height in both
transverse directions inside $[-8,8]^2$.

Consequently there are finitely many crossing families
$\mathcal F_1,\ldots,\mathcal F_J$ such that every macroscopic path
contains a member of one of them. Summing admissible weights gives
\[
 M(N)\leq J\sum_{j=1}^J\Mod_2(\mathcal F_j,T_N)
       \leq J^2\max_j\Mod_2(\mathcal F_j,T_N).
\]
One strip realizes this bound along a subsequence. Affine coordinates
give the asserted rectangles.
\end{proof}

The edge $y=0$ is the top of $\cQ$. The rectangles
$[0,1]\times[-1,0]$ and $[0,1]\times[1,2]$ are its top and bottom
bands, respectively. Coordinates serve only interpolation and topology:
all distances refer to $d$.

\begin{lemma}[Crossings and blockers]
\label{lem:planar-blocker}
In a closed coordinate rectangle:
\begin{enumerate}[label=\textup{(\roman*)}]
\item every continuum meeting both lateral edges meets every
top-to-bottom path;
\item every compact set meeting all top-to-bottom paths contains
a continuum meeting both lateral edges.
\end{enumerate}
\end{lemma}
\begin{proof}
For (i), suppose a lateral continuum and a top-to-bottom path are
disjoint. Enlarge the rectangle, extending the continuum horizontally
and the path vertically to the corresponding new edges. The extensions
remain disjoint: a new horizontal and vertical segment can meet only at
an original corner, already a common point of the original sets.
Approximate the extended path, within its positive clearance from the
continuum, by a polygonal path with interior in the enlarged rectangle.
After subdivision at intersections it contains a simple polygonal arc
between the top and bottom edges. The Jordan curve theorem, applied
after adjoining a boundary route, shows that this arc separates the
lateral edges. The connected continuum cannot avoid it.

For (ii), subdivide into finer square grids and mark the closed squares
meeting the blocker $K$. No edge-adjacent chain of unmarked squares
joins top to bottom. Label $A$ the unmarked squares reachable from the
top through unmarked edge-adjacent squares, and append a top row of
$A$ squares. Label all remaining squares and an appended bottom row $B$.
Draw the interfaces between unlike labels. Interior vertices have even
degree; each lateral boundary has an odd number of ends, since its
labels begin with $A$ and end with $B$. There are no top or bottom ends.
Some interface component therefore joins the lateral boundaries:
otherwise every component contributes an even number of ends to the
only lateral boundary it meets.

Along a path in this component, the adjacent $B$ squares are original
marked squares. An adjacent unmarked square would be reachable, while
an interface with the appended bottom row would give an unmarked
top-to-bottom chain. Successive $B$ squares meet at least at a corner.
Their union is a connected compact lateral crossing. A Hausdorff limit
as the mesh tends to zero is connected, meets both lateral edges, and
lies in $K$.
\end{proof}

The next finite-dimensional argument belongs to discrete extremal
length. Cannon--Floyd--Parry describe optimal weights through crossing
vectors and their pairing with transverse cuts
\cite[Sections 2.3 and 2.4.3]{CFPSquaring}; the probability formulation
of Albin--Poggi-Corradini gives the same least-norm incidence-vector
identity \cite[Theorem 5.1]{AlbinPoggiCorradini}. We give the argument
for the cell families used here.

\begin{lemma}[Crossing dual weights]
\label{lem:dual-weights}
For each retained $N$ there are nonnegative weights
$\rho_q=\rho_q^{(N)}$, $q\in T_N$, admissible for top-to-bottom
crossings of $\cQ$, such that
\begin{equation}
 \sum_{q\in T_N}\rho_q^2\leq \frac{C}{M(N)}.
 \label{eq:dual-energy}
\end{equation}
\end{lemma}
\begin{proof}
Let $P$ minimize Euclidean norm in the convex hull of the finitely many
cell-incidence vectors $v$ of left-to-right crossings. Its coordinates
are nonnegative and sum to at least one, so $P\ne0$. Minimality gives
$\langle P,v\rangle\geq\|P\|^2$, making $P/\|P\|^2$ admissible.
Conversely any admissible $\omega$ satisfies
$\langle\omega,P\rangle\geq1$. Thus the left-to-right modulus is
$\|P\|^{-2}$. By \Cref{lem:planar-blocker}(i), each top-to-bottom
incidence vector has inner product at least one with every $v$, hence
with $P$. Set $\rho=P$ and apply \Cref{lem:rectangle-localization}.
\end{proof}

\subsection{Continuous functions with separating levels}

For $\pi\subset\Qhat$ write
$\ell_\rho(\pi)=\sum_{q:B(q,2r_N)\cap\pi\ne\varnothing}\rho_q$.
For $x\in\cQ$, let $b_N(x)$ be the infimum of this cost over paths
from the top edge to $x$ in $\cQ$. Define the raw function
$\min\{1,b_N\}$, assigning exactly zero on the top edge, and extend
it by zero and one to the top and bottom bands. Admissibility gives
value one on the bottom edge.

Along every path in $\Qhat$, the difference of raw endpoint values is
at most its cost. Inside $\cQ$ this follows by concatenating with
nearly minimizing paths and clipping at one; a path beginning on the
top edge is itself a competitor. A path visiting both bands contains
a crossing of cost at least one. If it visits only the top band, each
endpoint value is bounded by the cost of its segment to that band.
If it visits only the bottom band, each endpoint value is at least one
minus that segment's cost, by concatenation with a top-to-endpoint
competitor. In either case the endpoint difference is bounded by the
cost of the entire path.

Triangulate $\Qhat$ by a finite coordinate grid containing the band
boundaries, with every triangle of $d$-diameter less than $r_N$.
Interpolate the raw vertex values affinely to obtain a continuous
$U_N^0:\Qhat\to[0,1]$. It is constant on the bands. Continuity alone
does not ensure that one selected level continuum meets every
straddling path. The following two path envelopes give the strict
sublevel and superlevel sets the connections to the bands needed for
that conclusion. Define
\[
 F_N(x)=\inf_{\pi:x\to\mathrm{top}}\max_\pi U_N^0,
 \qquad
 H_N(x)=\sup_{\pi:x\to\mathrm{bottom}}\inf_\pi F_N,
\]
where paths remain in $\Qhat$ and end in the indicated band.

\begin{lemma}[Regularization and level continua]
\label{lem:level-continua}
The functions $F_N,H_N$ are continuous, with the prescribed band values,
and for every path $K\subset\Qhat$,
\begin{equation}
 \osc_K H_N\leq\osc_K F_N\leq\osc_K U_N^0.
 \label{eq:flood-osc}
\end{equation}
For every $0<t<1$, the set $\{H_N=t\}\cap\cQ$ contains a lateral
crossing continuum $C_{N,t}$ of diameter at least a fixed $d_1>0$.
Every path in $\Qhat$ whose endpoint values strictly straddle $t$
meets $C_{N,t}$.
\end{lemma}
\begin{proof}
We have $F_N\geq U_N^0$ and $H_N\leq F_N$; constant paths give
the band values. Concatenation with nearly optimizing paths gives,
for $x,y\in K$,
\[
 F_N(x)\leq\max\{\max_K U_N^0,F_N(y)\},\qquad
 H_N(y)\geq\min\{\inf_K F_N,H_N(x)\}.
\]
If $F_N(x)>F_N(y)$, the first inequality places their values between
$\min_K U_N^0$ and $\max_K U_N^0$. If $H_N(x)>H_N(y)$, the
second places them between $\inf_K F_N$ and $\sup_K F_N$.
This proves \eqref{eq:flood-osc}. Short coordinate segments and
continuity of $U_N^0$ then prove continuity of both functions,
including at the boundary.

Every point of $\{H_N>t\}$ joins the bottom band within that set:
choose a path with $\inf F_N>t$ and use its tails. Likewise every
point of $\{F_N<t\}$ joins the top within $\{F_N<t\}$.
To connect a point $x$ with $H_N(x)<t\leq F_N(x)$ to the top,
choose $H_N(x)<s<t$ and follow any top-going path until its first
point with $F_N=s$. On this prefix $F_N\geq s$, so concatenation
with nearly maximizing paths from any prefix point $z$ gives
\[
 H_N(x)\geq\min\{s,H_N(z)\}.
\]
Hence $H_N(z)\leq H_N(x)<t$ throughout the prefix. Its endpoint
joins the top through $F_N<t$, hence through $H_N<t$. Thus both
strict sides of the level join their respective bands, and then the
outer top and bottom edges, without meeting the level.

The compact set $\{H_N=t\}$ blocks all top-to-bottom paths by
continuity and lies in $\cQ$. By \Cref{lem:planar-blocker}(ii) it
contains a lateral continuum $C_{N,t}$; take $d_1$ to be the distance
between the lateral edges of $\cQ$. Complete any straddling path to
a top-to-bottom path using the strict-side connections. The completed
path meets $C_{N,t}$ by \Cref{lem:planar-blocker}(i), and the added
pieces avoid the level, proving the last assertion.
\end{proof}

Figure~\ref{fig:level-continuum} isolates the last step: the strict-side
extensions cannot account for the intersection with the selected
continuum, so the original straddling path must meet it.

\begin{figure}[htbp]
\centering
\begin{tikzpicture}[x=1cm,y=1cm,font=\small,
  line cap=round,line join=round,>=Stealth]
  \definecolor{levelblue}{RGB}{54,101,125}
  \definecolor{levelochre}{RGB}{164,111,48}
  \fill[black!3] (0,1.25) rectangle (7,1.72);
  \fill[black!3] (0,-1.72) rectangle (7,-1.25);
  \draw[black!45] (0,-1.72) rectangle (7,1.72);
  \draw[black!45] (0,1.25)--(7,1.25) (0,-1.25)--(7,-1.25);
  \node at (1.9,1.49) {top band: $H_N=0$};
  \node at (1.9,-1.49) {bottom band: $H_N=1$};
  \node[anchor=north west,text=black!65] at (.10,1.15) {$\mathcal Q$};
  \draw[levelblue,line width=2pt]
    (0,.13)--(.37,.03)--(.69,.19)--(1.04,.05)--(1.34,.20)
    --(1.64,-.04)--(1.98,.08)--(2.28,-.09)--(2.60,.04)
    --(2.91,-.13)--(3.19,.06)--(3.53,-.09)--(3.84,.09)
    --(4.16,-.02)--(4.48,.15)--(4.81,.01)--(5.13,.18)
    --(5.47,.04)--(5.80,.20)--(6.10,.05)--(6.41,.20)
    --(6.70,.08)--(7,.20);
  \draw[levelblue,line width=1.1pt]
    (1.34,.20)--(1.42,.37)--(1.59,.32)
    (2.28,-.09)--(2.18,-.28)--(2.38,-.24)
    (4.81,.01)--(4.89,-.21)--(5.03,-.15)
    (6.10,.05)--(6.04,-.14)--(6.22,-.19);
  \node[levelblue,anchor=south] at (1.95,.40) {$C_{N,t}\subset\{H_N=t\}$};
  \draw[levelochre,dashed,thick]
    (4.9,1.72) .. controls (4.4,1.32) and (3.88,1.29) .. (3.72,.79);
  \draw[levelochre,thick]
    (3.72,.79) .. controls (3.50,.34) and (4.24,-.24) .. (4.10,-.80);
  \draw[levelochre,dashed,thick]
    (4.10,-.80) .. controls (4.00,-1.24) and (4.94,-1.29) .. (5.1,-1.72);
  \fill[levelochre] (3.72,.79) circle (1.5pt) (4.10,-.80) circle (1.5pt);
  \node[anchor=west] at (3.96,.74) {$H_N<t$};
  \node[anchor=east] at (3.89,-.86) {$H_N>t$};
  \node[levelochre,anchor=west] at (4.27,-.46) {$\eta$};
  \node[anchor=west,text=black!85,align=left] at (7.32,.82)
    {strict-side\\extension};
  \draw[black!65,thin] (7.2,.86)--(5.10,1.18);
  \node[anchor=west,text=black!85,align=left] at (7.32,-.85)
    {strict-side\\extension};
  \draw[black!65,thin] (7.2,-.82)--(5.05,-1.36);
\end{tikzpicture}
\caption{A separating level continuum and a straddling path. The dashed
extensions stay in the strict sublevel and superlevel sets and complete
$\eta$ to a top-to-bottom crossing. That crossing meets $C_{N,t}$, so
$\eta$ itself meets $C_{N,t}$. The continuum is drawn schematically;
neither smoothness nor path connectedness is assumed.}
\label{fig:level-continuum}
\end{figure}


\subsection{Oscillation and annular estimates}

The level continua now let us estimate oscillation by crossing an annulus.
Energies formed from image diameters and their relation to condenser
capacity appear in Pansu's coarse conformal geometry \cite[\S6]{Pansu1989};
the estimates below use the finite-cell energy of our regularized scalar
functions. Put $U=\operatorname{int}_Z\Qhat$. For a continuous function $f$ on
an open set containing a compact set $K$, and $J\geq2$, define
\[
 \mathcal E_n^{(J)}(f;K)=
 \sum_{\substack{q\in T_n\\B(q,2r_n)\cap K\ne\varnothing}}
       \bigl(\osc_{\overline B(q,Jr_n)}f\bigr)^2.
\]
The closed balls lie in the domain for all sufficiently large $n$.

\begin{lemma}[Local oscillation]
\label{lem:local-oscillation}
For fixed compact $K\subset U$ and $J>0$, uniformly for $x\in K$
and sufficiently large $N$,
\begin{equation}
 \osc_{\overline B(x,Jr_N)}H_N
 \leq\sum_{\substack{q\in T_N\\d(q,x)\leq C(J)r_N}}\rho_q.
 \label{eq:local-osc}
\end{equation}
Here $C(J)$ depends only on $J$ and the ball path constant.
Consequently, for every fixed $g\in G$, compact $K'\subset gU$,
and $J\geq2$,
\[
 \mathcal E_N^{(J)}(H_N\circ g^{-1};K')
       \leq\frac{C(g,K',J)}{M(N)}\longrightarrow0.
\]
\end{lemma}
\begin{proof}
Join two points of $\overline B(x,Jr_N)$ inside $B(x,LJr_N)$
by \Cref{lem:paths}. Vertices of triangles meeting this path lie
in $B(x,(LJ+1)r_N)$; any two such vertices join inside
$B(x,L(LJ+1)r_N)\subset U$ for large $N$. The raw endpoint
inequality bounds every difference of their values by the right side
of \eqref{eq:local-osc}, with $C(J)=L(LJ+1)+3$. It therefore bounds
the interpolant's oscillation along the original path. Apply
\eqref{eq:flood-osc}.

Squaring and summing costs a bounded factor by \Cref{lem:packing}:
both the number of nearby weights and the number of centers charging
each weight are bounded. The same holds for $br_N$-separated centers,
with constants depending on $b>0$. A fixed $g$ is bi-Lipschitz,
so inverse images of target balls lie in source balls of radius
$C_gJr_N$, their centers are $b_gr_N$-separated, and all these source
balls lie in a fixed compact subset of $U$. The summation bound and
\eqref{eq:dual-energy} finish the proof.
\end{proof}

The mesh-scale estimate preserves the small dual energy. We next
compare oscillation across a ball of radius $R=r_l$ with the dual
energy in a surrounding annulus. Expansion takes that annulus to a
fixed-size region and takes the mesh $r_N$ to $r_{N-l}$; this is why
the next bound contains $M(N-l)$.

\begin{lemma}[Annular estimate]
\label{lem:annular}
For every compact $K\subset U$ there are $R_*>0$ and an integer
$m_0$ such that, if $x\in K$, $R=r_l\leq R_*$, and $N-l\geq m_0$,
then any two points joined by a path in $B(x,R/2)$ have
$H_N$-value difference $\Delta$ satisfying
\begin{equation}
 \Delta^2\leq C M(N-l)
       \sum_{\substack{q\in T_N\\R/4\leq d(q,x)\leq4R}}\rho_q^2.
 \label{eq:annular}
\end{equation}
The constant is uniform; $N$ is sufficiently large for the fixed
interior region.
\end{lemma}
\begin{proof}
Choose $R_*$ with $4R_*<d_1$ and all closed $5R_*$-balls about
$K$ contained in $U$. All following restrictions on $m_0$ concern
only $r_N/R=a^{-(N-l)}$.

Expand at $x,l$ by $h$, and choose macroscopic admissible weights
$\omega$ at scale $N-l$ with squared mass $M(N-l)$. Put
\[
 \mathcal A=\{q\in T_N:R/2\leq d(q,x)\leq3R\},\qquad
 \sigma_q=\sum_{\substack{\xi\in T_{N-l}\\
               d(\xi,hq)\leq C_0r_{N-l}}}\omega_\xi,
 \quad q\in\mathcal A,
\]
where $C_0>C_E+2$. Expansion on $B(x,3R)$ separates the $hq$ by
a fixed multiple of $r_{N-l}$. Packing bounds both overlap
multiplicities, giving
\[
 \sum_{q\in\mathcal A}\sigma_q^2\leq C M(N-l).
\]

Let $I$ be the open interval between the endpoint values. For $t\in I$,
the continuum $C_{N,t}$ meets the given path and exits $B(x,2R)$,
since its diameter exceeds $4R$. By \Cref{lem:paths}, join these
two continuum points in its $r_N/4$-neighborhood and retain a radial
subpath $\eta_t$ between radii $R$ and $2R$, contained in that closed
annulus. Its diameter is at least $R$, so $h\eta_t$ is macroscopic
by the choice of $d_0$ in \Cref{sec:geometry}.

If a target cell $B(\xi,2r_{N-l})$ meets $h\eta_t$ at $hy$,
choose $q\in T_N$ with $d(q,y)\leq r_N$. For large $m_0$, this
center belongs to $\mathcal A$ and
$d(hq,hy)\leq C_Er_N/R$, so $\omega_\xi$ occurs in $\sigma_q$.
Some level point lies within $r_N/4$ of $y$, hence in
$\overline B(q,2r_N)$. Thus $t$ lies in the closed range interval
$I_q$ of $H_N$ on that ball. Accounting for all target cells gives
\[
 1\leq\sum_{q\in\mathcal A}\sigma_q\mathbf1_{I_q}(t)
 \quad(t\in I),
 \qquad
 \Delta\leq C M(N-l)^{1/2}
       \left(\sum_{q\in\mathcal A}|I_q|^2\right)^{1/2}.
\]
The second inequality follows by integration and Cauchy--Schwarz.
Apply \eqref{eq:local-osc} with $J=2$ and its bounded-overlap
summation. Taking $C(2)r_N<R/4$ confines every charged center to
$R/4\leq d(q,x)\leq4R$, proving \eqref{eq:annular}.
\end{proof}

\subsection{The limiting function and measure}

The finite measures
$\mu_N=M(N)\sum_{q\in T_N}\rho_q^2\delta_q$
have uniformly bounded mass by \eqref{eq:dual-energy}.

\Needspace{5\baselineskip}
\begin{proposition}[Scalar limit]
\label{prop:scalar-limit}
Along a subsequence of the record indices,
$H_N\to u$ locally uniformly on $U$ and $\mu_N\rightharpoonup\mu$,
where $u:U\to[0,1]$ is continuous and nonconstant and $\mu$ is finite.
The function $u$ is zero and one on the interiors of the respective
bands, and is locally H\"older continuous if $\lambda>1$.
For every compact $K\subset U$ and sufficiently small
$s>0$, uniformly for $x\in K$,
\begin{equation}
 \bigl(\osc_{\overline B(x,s)}u\bigr)^2
       \leq C\mu\bigl(\overline B(x,C_\mu s)\bigr).
 \label{eq:measure}
\end{equation}
For each $0<t<1$ there is a continuum $C_t\subset\cQ$ of diameter
at least $d_1$, with $u=t$ on $C_t\cap U$, meeting every compact
path in $U$ whose endpoint values strictly straddle $t$.
The fixed-transform energy bounds of \Cref{lem:local-oscillation}
hold along the retained subsequence.
\end{proposition}
\begin{proof}
Fix compact $K\subset U$, and let $x,y\in K$ with $d(x,y)\leq r_k$.
They join inside $B(x,Lr_k)$. Choose $b$ with $a^b>2L$ and consider
indices $l\leq k-b$ above a fixed interior threshold, in an arithmetic
progression of step $b_1$ with $a^{b_1}>16$. The annuli
$\{r_l/4\leq d(x,\cdot)\leq4r_l\}$ are disjoint. For $N$
sufficiently large relative to $k$, \eqref{eq:annular} and
\eqref{eq:records} give
\[
 |H_N(x)-H_N(y)|^2\leq C\lambda^{-l}
 \mu_N\{r_l/4\leq d(x,\cdot)\leq4r_l\}.
\]
Summing over at least $ck$ indices yields
$|H_N(x)-H_N(y)|\leq C_Kk^{-1/2}$. If $\lambda>1$, the single
choice $l=k-b$ gives $C_K\lambda^{-k/2}$ instead.

For each tolerance this controls a tail uniformly; the finitely many
omitted functions are individually uniformly continuous. Thus the
sequence is equicontinuous on $K$. Arzel\`a--Ascoli on a compact
exhaustion of $U$, followed by weak compactness of the measures on
$Z$, gives the asserted subsequence. The two constant band values
prove nonconstancy. The stronger estimate gives local H\"older exponent
$\log\lambda/(2\log a)$.

For \eqref{eq:measure}, fix $A>2L$ and choose
$As\leq R=r_l<aAs$. Any two points of $\overline B(x,s)$ join in
$B(x,Ls)\subset B(x,R/2)$. The annular and record estimates give
\[
 \bigl(\osc_{\overline B(x,s)}H_N\bigr)^2
       \leq C\mu_N\bigl(\overline B(x,4aAs)\bigr)
\]
for large $N$. Local uniform convergence passes the oscillation to
the limit; the closed-set inequality for weak convergence bounds the
right-hand limsup by $C\mu(\overline B(x,4aAs))$.
Take $C_\mu=4aA$.

For each fixed $t$, take a Hausdorff limit $C_t$ of $C_{N,t}$ in
$\cQ$. It is connected and has diameter at least $d_1$. Local
uniform convergence gives $u=t$ on $C_t\cap U$. A fixed compact
straddling path also straddles $t$ for every sufficiently large
$H_N$, hence meets those $C_{N,t}$ and their limit. The subsequence
used here may depend on $t$; no simultaneous selection is needed.
\end{proof}

\subsection{Two comparison tools}

The scalar construction will be used in three ways. The large continua
$C_t$ provide levels for annular constructions, with $u=t$ on
$C_t\cap U$; no value of $u$ is asserted outside $U$. The measure
estimate \eqref{eq:measure} turns oscillation on a small ball into a
lower bound for its enlarged ball's $\mu$-mass. To compare
transformations, we retain $H_N$ and $C_{N,t}$: the identity $H_N=t$
holds on the entire continuum, and the local oscillation energies
tend to zero after each transformation has been fixed. Those
comparisons will pass to subcontinua contained in compact subsets of
the common open domains.

The first tool integrates variation of a second function over a
family of level continua. It is a finite-cell statement, independent
of the preceding construction.

\begin{lemma}[A finite-cell level comparison]
\label{lem:cross-energy}
Let $f,w$ be continuous on an open set $V\subset Z$, let $K\subset V$
be compact, and let $I\subset\RR$ be an interval. Suppose that for
every $t\in I$ there is a continuum $D_t\subset K$ with $f=t$ on
$D_t$ and $\osc_{D_t}w\geq b_0>0$. For sufficiently large $n$,
\begin{equation}
 b_0|I|\leq
 \sum_{\substack{q\in T_n\\B(q,2r_n)\cap K\ne\varnothing}}
 \osc_{\overline B(q,2r_n)}f\,
 \osc_{\overline B(q,2r_n)}w.
 \label{eq:cross-energy}
\end{equation}
Consequently these hypotheses cannot hold for sequences with both
local oscillation energies on $K$ tending to zero, if $b_0$ and
$|I|$ stay bounded below by positive constants.
\end{lemma}
\begin{proof}
Let $I_q(f)$ be the closed range interval of $f$ on the closed cell.
For large $n$ these cells lie in $V$. The interval $w(D_t)$ is covered
by the range intervals of $w$ on cells meeting $D_t$, and each such
cell has $t\in I_q(f)$. Therefore
\[
 b_0\leq\sum_{q:B(q,2r_n)\cap K\ne\varnothing}
 \osc_{\overline B(q,2r_n)}w\,\mathbf1_{I_q(f)}(t).
\]
Integrate over $I$ and use
$|I\cap I_q(f)|\leq\osc_{\overline B(q,2r_n)}f$.
Only finitely many interval indicators are integrated, so no
measurable selection of $D_t$ is required. Cauchy--Schwarz gives
the final assertion.
\end{proof}

\begin{lemma}[Stopping a continuum at a sphere]
\label{lem:stop-continuum}
Let $D$ be a continuum, $p\in Z$, and $r>0$. If
$x\in D\setminus\overline B(p,r)$ and $D\cap B(p,r)\ne\varnothing$,
then $D$ contains a subcontinuum through $x$ meeting the radius-$r$
sphere and avoiding $B(p,r)$.
\end{lemma}
\begin{proof}
In a compact metric space $X$, a point's component is the intersection
$Q_x$ of its clopen neighborhoods. To verify connectedness of $Q_x$,
suppose it is separated into two compact pieces in disjoint open sets
$O_1,O_2$. Compactness of $X\setminus(O_1\cup O_2)$ gives a finite
intersection $W$ of clopen neighborhoods of $x$ contained in their
union. Its intersection with the $O_i$ containing $x$ is clopen in
$X$ and misses the other piece of $Q_x$, a contradiction. The same
compactness argument shows that every open neighborhood of a component
contains a clopen neighborhood of it.

Apply this to $X=D\setminus B(p,r)$ and the component $C$ of $x$.
If $C$ missed the sphere, it would have a clopen neighborhood $W$ in
$X$ missing the sphere. Compactness places $W$ a positive distance
from $\overline B(p,r)$, so $W$ is also clopen in $D$. It is nonempty
and misses $D\cap B(p,r)$, contradicting connectedness. Thus $C$ is
the required subcontinuum.
\end{proof}
